\documentclass[12pt]{scrartcl}

\usepackage[utf8]{inputenc}
\usepackage[T1]{fontenc}
\usepackage{lmodern}
\usepackage{graphicx}
\usepackage[pdftex,hyperref,dvipsnames]{xcolor}
\usepackage{listings}
\usepackage[a4paper,lmargin={2cm},rmargin={2cm},tmargin={3.5cm},bmargin = {2.5cm},headheight = {4cm}]{geometry}
\usepackage{amsmath,amssymb,amstext,amsthm}
% \usepackage[lined,algonl,boxed]{algorithm2e} 
\usepackage{algpseudocode}
\usepackage{tikz}
\usepackage{pgfplots}
\usepackage{hyperref}
\usepackage{url}
\usepackage{tcolorbox}
\usepackage{float}
\usepackage[inline]{enumitem}
\usepackage[headsepline]{scrlayer-scrpage} 

\usetikzlibrary{automata,positioning}

\usepgfplotslibrary{external}
%\tikzexternalize
\pgfplotsset{width=10cm,compat=1.9}
\pgfplotsset{every axis/.style={
    axis line style = thick,
    grid = both,
    tick pos = left,
    tick align = outside,
    tick style = {thick,black},
    fill = blue,
	line width = 2pt,
}}

\lstset{
  basicstyle=\ttfamily\small,
  numbers=left,
  numberstyle=\tiny,
  stepnumber=1,
  numbersep=8pt,
  frame=single,
  breaklines=true,
  columns=fullflexible
}

\lstset{
  language=C,
  basicstyle=\ttfamily\small,
  keywordstyle=\color{blue},
  commentstyle=\color{gray},
  stringstyle=\color{green!60!black},
  numbers=left,
  numberstyle=\tiny,
  stepnumber=1,
  numbersep=8pt,
  frame=single,
  breaklines=true,
  columns=fullflexible
}

\renewcommand{\theenumi}{(\alph{enumi})}
\renewcommand{\labelenumi}{\text{\theenumi}}
\renewcommand{\labelenumii}{(\roman{enumii})} 

\ohead{Weiting Wang (3855168)\\
       Nico Strohbach (3546914)\\
       Jan Theil (3794698)}

\ihead{Reinforcement Learning}

\newcounter{exnum}
\newcommand{\exercise}[1]{\section*{Problem \theexnum\stepcounter{exnum}: #1}}

\begin{document}

\setcounter{exnum}{1}
\exercise{Planning and Learning}

\begin{enumerate}
  \item In Dyna-Q+ an exploration bonus is added to states during planning that have not been visited in a long time.
        Therefor, it tends to visit and explore states, that have not been visited in a long time.
        In the first phase, this leads to a faster "accurate" model since it explores all states
        more broadly compared to Dyna-Q. Hence, it finds the optimal solution faster.
        This is why the slope of the red curve goes up sooner than the slope of the blue curve
        but the differences of these curves stays somewhat stationary.
        In the second phase a new optimal path is possible due to the openning in the right side.
        Since Dyna-Q+ keeps exploring the right side more often than Dyna-Q (because of the exploration bonus)
        it finds the new optimal path faster while Dyna-Q tends to stick to a once found optimal path.
  \item Instead of storing a single $R, S'$ entry in the Model table in the line $Model(S, A) \leftarrow R, S'$ we try
        to store an estimation of the probability distribution of the stochastic environment.
        This could be done by e.g. keeping counters of each possible $R, S'$-pair and updating them accordingly.
        So, instead of a single entry we would have another (sub-)row with which an estimation of $p(\cdot)$ could be calculated
        by $p(\cdot) \approx \dfrac{N(R, S')}{\sum_{R, S} N(R, S)}$. In the planning phase, this estimation is used to obtain new $R, S'$ pairs.
        This modification would not perform well on changing environments since the once learned probability distribution estimates
        would change very slowly when new distributions are introduced in the environment. This is the same problem as with normal Dyna-Q learning.
        To perform well on both stochastic and changing environments, a combination of the modification above and the Dyna-Q+ semantics
        would be useful. Therefor, something like an exploration bonus would help with changing environments.
\end{enumerate}

\exercise{Monte Carlo Tree Search on the Taxi environment}

\end{document}